\section{Problem Definition}

\subsection{Problem Statement}
Small and medium distributors and multi-SKU retailers must decide, for every product and store, how much stock to reorder and when. Two costly errors follow from getting this wrong. Ordering too much ties up cash and leads to storage costs, markdowns, and spoilage. Ordering too little causes stockouts, lost sales, and unhappy customers. Because purchasing must happen before demand is known, the decision depends on a forecast of what will sell during the supplier lead time plus the review period.

In practice, many businesses make this decision with spreadsheets, experience, or a simple average of recent sales. These methods have three weaknesses. First, they ignore the factors that actually drive demand, such as promotions, holidays, day-of-week patterns, and seasonality. Second, they produce a single number with no measure of uncertainty, so the buffer stock added on top is a guess rather than something derived from forecast error. Third, they do not scale: a retailer with thousands of store--item combinations cannot tune each rule by hand.

This project addresses the question: \emph{what will likely sell in the next few days or weeks, how uncertain is that, and how much should I reorder?} We build a web application that learns from historical sales data, forecasts near-future demand with an uncertainty range (P10, P50, and P90), and converts the forecast into a reorder recommendation using safety stock calculated from validation error. The system supports the manager's decision; it does not place orders automatically and does not replace a POS or ERP system.

\subsection{Target Audience}
The primary users are inventory and purchasing managers at small and medium distributors and multi-SKU retail businesses who currently plan orders manually and lack a dedicated data science team. A secondary audience is business analysts who need transparent forecast accuracy figures to justify purchasing decisions. The tool takes a CSV export of historical sales as input, so it can be used alongside any existing POS or ERP system without integration work.

\subsection{AI Justification}
A traditional rule-based approach, such as a fixed reorder point or a moving average of recent sales, is simple but is not sufficient for this problem, for the following reasons.

\begin{itemize}
    \item \textbf{Demand depends on many interacting drivers.} Sales vary with promotions, school and public holidays, day of week, month, store characteristics, and recent trend, and these effects interact. For example, a promotion has a different effect on a Saturday than on a Monday. Encoding all such interactions as hand-written rules is impractical, whereas gradient-boosted trees such as LightGBM learn them directly from data.
    \item \textbf{Scale and shared patterns.} A retailer may have thousands of store--item series. A single global model trained on all series can transfer patterns between similar series, which helps those with short or noisy histories, while rules would need to be tuned separately for each one.
    \item \textbf{Uncertainty is needed, not just a point forecast.} Safety stock and stockout alerts depend on how wrong the forecast might be. Quantile regression produces P10, P50, and P90 forecasts, and the interval coverage can be checked against the target of about 80\%. Fixed rules provide no principled estimate of this uncertainty.
    \item \textbf{Adaptation without manual re-engineering.} When demand patterns change, the model is retrained on new data under time-based validation. A rule set would need to be revised by hand.
    \item \textbf{The benefit can be measured.} We keep the rule-style methods as baselines (seasonal naive and moving average) and compare LightGBM against them using MAE, RMSE, WMAPE/MASE, and interval coverage on a time-based holdout. The case for machine learning therefore rests on measured improvement and not on assumption.
\end{itemize}

The forecast is then turned into a decision using standard inventory logic. With lead time $L$ and review period $R$, let $T = L + R$ and let $D_T$ be the P50 demand over $T$ days. Safety stock is $SS = z\,\sigma_d\sqrt{T}$, where $\sigma_d$ is the daily validation RMSE and $z$ is the service-level factor. The recommended order quantity is $Q = \max(0, \lceil D_T + SS - S - I \rceil)$, where $S$ is stock on hand and $I$ is incoming stock. Here the machine learning model supplies the demand forecast and the error estimate, and the reorder rule remains transparent and auditable.

\subsection{Scope and Limitations}
The project uses two public Kaggle datasets: Rossmann Store Sales (store-level daily sales with promotion, holiday, and store information) and Store Item Demand Forecasting (store--item daily sales). Neither dataset contains inventory levels or lead times, so these values are simulated or user-entered and are labeled as such in the application. Neither dataset represents Pakistani retail, and we make no such claim. Evaluation uses time-based splits only, with no random splitting, to avoid leakage from future data.
